\chapter{The Remaining Formation Problem}
\label{ch:outlook}

The constructions describe agents that compare laws, receive internal records, and inherit descriptions across scales. A normalized predictor makes model agreement operational; an inherited update makes a retained description dynamical. Neither follows from identifying a coherent block at one instant.

The quantifiers in \cref{inh:recursive-formation} distinguish a prescribed partition from a selected one, a transient block from a persistent one, and agreement of one record law from agreement of subsequent responses. Likewise, a smaller model-sector mobility imposes a timescale separation without explaining its origin.

\Cref{ch:fusion} supplies an exact two-stage example after a change of interaction: square-root KL costs produce a nonsmooth force and finite-time equality. Local models, parent Fisher metrics, contact multiplicities, and fresh-response kernels are inherited. The initial class is restricted, and all agents eventually fuse. This mechanism does not establish formation under the smooth potential of \cref{ch:dynamics}.

\Cref{ch:smooth-closure} gives exact general-law reductions under a specified conditional reconstruction architecture. Conditional log-density spectra retain the hidden information entering the marginal force; record marks complete the state for declared feedback protocols. Equal spectra permit different internal conditionals and correlated laws while public marginals follow the original coarse Fisher flow. That equality is an invariant preparation and is not first reached at finite time.

\Cref{ch:smooth-erasure} instead lets an internal gate enter a prescribed flat channel region, after which private laws can be discarded. Source preferences change into persistent pairs whose public states remain distinct at finite times. This first-stage reduction uses a finite product family and no fusion penalty. It supplies the erasing architecture, retains every public child role, and excludes later private-sensitive feedback. Its two invariant public regions do not merge again.

The unresolved collective problem is aggregation of distinct public states under interfaces selected by the dynamics, with inherited responses and updates that permit another formation step. Fine states with identical proposed parent data must have identical retained derivatives and responses. The arithmetic-mixture counterexample shows why this can fail even for identical child roles. Missing information may concern conditional dependence, interface response, or shared-source memory; it need not fit a predetermined distribution family.

Deriving a base geometry also remains separate. State-dependent interactions at one context do not specify compatible transports between nearby contexts. A curved connection requires that additional structure; a common geometry across agents requires compatibility between their structures.
